% Options for packages loaded elsewhere
\PassOptionsToPackage{unicode}{hyperref}
\PassOptionsToPackage{hyphens}{url}
\documentclass[
  ignorenonframetext,
]{beamer}
\newif\ifbibliography
\usepackage{pgfpages}
\setbeamertemplate{caption}[numbered]
\setbeamertemplate{caption label separator}{: }
\setbeamercolor{caption name}{fg=normal text.fg}
\beamertemplatenavigationsymbolsempty
% remove section numbering
\setbeamertemplate{part page}{
  \centering
  \begin{beamercolorbox}[sep=16pt,center]{part title}
    \usebeamerfont{part title}\insertpart\par
  \end{beamercolorbox}
}
\setbeamertemplate{section page}{
  \centering
  \begin{beamercolorbox}[sep=12pt,center]{section title}
    \usebeamerfont{section title}\insertsection\par
  \end{beamercolorbox}
}
\setbeamertemplate{subsection page}{
  \centering
  \begin{beamercolorbox}[sep=8pt,center]{subsection title}
    \usebeamerfont{subsection title}\insertsubsection\par
  \end{beamercolorbox}
}
% Prevent slide breaks in the middle of a paragraph
\widowpenalties 1 10000
\raggedbottom
\AtBeginPart{
  \frame{\partpage}
}
\AtBeginSection{
  \ifbibliography
  \else
    \frame{\sectionpage}
  \fi
}
\AtBeginSubsection{
  \frame{\subsectionpage}
}
\usepackage{iftex}
\ifPDFTeX
  \usepackage[T1]{fontenc}
  \usepackage[utf8]{inputenc}
  \usepackage{textcomp} % provide euro and other symbols
\else % if luatex or xetex
  \usepackage{unicode-math} % this also loads fontspec
  \defaultfontfeatures{Scale=MatchLowercase}
  \defaultfontfeatures[\rmfamily]{Ligatures=TeX,Scale=1}
\fi
\usepackage{lmodern}
\ifPDFTeX\else
  % xetex/luatex font selection
\fi
% Use upquote if available, for straight quotes in verbatim environments
\IfFileExists{upquote.sty}{\usepackage{upquote}}{}
\IfFileExists{microtype.sty}{% use microtype if available
  \usepackage[]{microtype}
  \UseMicrotypeSet[protrusion]{basicmath} % disable protrusion for tt fonts
}{}
\makeatletter
\@ifundefined{KOMAClassName}{% if non-KOMA class
  \IfFileExists{parskip.sty}{%
    \usepackage{parskip}
  }{% else
    \setlength{\parindent}{0pt}
    \setlength{\parskip}{6pt plus 2pt minus 1pt}}
}{% if KOMA class
  \KOMAoptions{parskip=half}}
\makeatother
\setlength{\emergencystretch}{3em} % prevent overfull lines
\providecommand{\tightlist}{%
  \setlength{\itemsep}{0pt}\setlength{\parskip}{0pt}}
% Slide layout defaults for warning-clean scientific decks.
\usepackage{etoolbox}
\IfFileExists{xurl.sty}{\usepackage{xurl}}{}
\IfFileExists{seqsplit.sty}{\usepackage{seqsplit}}{\newcommand{\seqsplit}[1]{#1}}
\protected\def\breakseq#1{\seqsplit{#1}}
\protected\def\breaktt#1{\begingroup\ttfamily\seqsplit{#1}\endgroup}
\setlength{\emergencystretch}{6em}
\tolerance=5000
\hbadness=10000
\hfuzz=1pt
\setlength{\tabcolsep}{2pt}
\AtBeginEnvironment{longtable}{\tiny\renewcommand{\arraystretch}{0.86}\setlength{\tabcolsep}{1pt}}
\AtBeginEnvironment{tabular}{\tiny\renewcommand{\arraystretch}{0.86}\setlength{\tabcolsep}{1pt}}

% Natbib and cross-reference fallbacks — slides don't load natbib
% or cleveref, but combined-PDF manuscript prose may emit these
% commands. The fallback renders citations as a bracketed key list
% and cross-references as detokenized labels so slides don't fail on
% undefined control sequences or raw underscores. \providecommand is
% a no-op if packages load later.
\providecommand{\citep}[1]{[#1]}
\providecommand{\citet}[1]{#1}
\providecommand{\citealp}[1]{#1}
\providecommand{\citeauthor}[1]{#1}
\providecommand{\citeyear}[1]{#1}
\providecommand{\cref}[1]{\texttt{\detokenize{#1}}}
\providecommand{\Cref}[1]{\texttt{\detokenize{#1}}}
\usepackage{bookmark}
\IfFileExists{xurl.sty}{\usepackage{xurl}}{} % add URL line breaks if available
\urlstyle{same}
% fallback for those not using the hyperref driver hyperxmp:
\makeatletter
\@ifundefined{xmpquote}{\newcommand{\xmpquote}[1]{#1}}{}
\makeatother
\hypersetup{
  hidelinks,
  pdfcreator={LaTeX via pandoc}}

\author{\texorpdfstring{}{}}
\date{}

\begin{document}

\section{Full Text, Language, and
Embeddings}\label{full-text-language-and-embeddings}

\begin{frame}[allowframebreaks]{Full Text, Language, and Embeddings}
Beyond bibliographic metadata, the pipeline mines the textual content of
each record. This stage bridges the gap between a bibliographic
inventory and a semantic understanding of the literature.
\end{frame}

\begin{frame}[fragile,allowframebreaks]{Full-Text Availability}
\protect\phantomsection\label{full-text-availability}
An open-access resolver maps each record to a downloadable PDF where one
exists (a known \texttt{pdf\_url}, or an Unpaywall lookup by DOI), and
an opt-in, network-gated downloader fetches it to a deterministic path.
Full-text availability is summarized without requiring any download, so
the offline default still reports coverage. For this run:

\begin{itemize}
\tightlist
\item
  \textbf{Abstract coverage}: 55.5\% of records (1277 of

  \begin{enumerate}
  [1)]
  \setcounter{enumi}{2301}
  \tightlist
  \item
    carry an abstract; 1025 records lack one.
  \end{enumerate}
\item
  \textbf{Open-access status}: 14.4\% of records are open access (331
  records); the remainder are closed or unknown.
\item
  \textbf{PDF availability}: 40.9\% of records (941) have a direct PDF
  link; 940 have a publisher PDF, and 1361 have no full-text source
  available.
\end{itemize}

The identifier coverage for this corpus is: 2248 DOIs, 932 OpenAlex IDs,
and 1 arXiv IDs. DOI coverage dominates, enabling robust cross-engine
de-duplication.
\end{frame}

\begin{frame}[allowframebreaks]{Language and Entity Extraction}
\protect\phantomsection\label{language-and-entity-extraction}
Titles, abstracts, and (when present) full text are tokenized and
reduced to keyphrases and named entities by offline, dependency-light
extractors --- no mandatory LLM. Term-frequency statistics drive a
TF-IDF representation over a 500-feature vocabulary. The most frequent
terms in the corpus are: modafinil, treatment, study, effects, patients,
results, sleep, used, use, drug, studies, clinical, mg, using, placebo,
cognitive, associated, effect, however, disorder. These terms reflect
the clinical, pharmacological, and cognitive vocabulary characteristic
of the modafinil literature.
\end{frame}

\begin{frame}[fragile,allowframebreaks]{Embeddings}
\protect\phantomsection\label{embeddings}
Every title, abstract, and full text is embedded into a shared vector
space by a deterministic, offline method --- TF-IDF followed by
truncated SVD, i.e.~latent semantic analysis
\breakseq{[@deerwester1990indexing]}. The embedding dimensionality is 50
components (configurable via
\breaktt{project\_config.embeddings.n\_components}), and the TF-IDF
vocabulary is capped at 500 features (configurable via
\breaktt{project\_config.embeddings.max\_features}). The embedding is
byte-stable across runs: the same input text always yields identical
vectors, so the derived similarity matrix, nearest-neighbour lists,
clusters, and two-dimensional projection are all reproducible.

An optional transformer backend can be enabled by setting
\breaktt{project\_config.embeddings.method:\ transformer} (requires the
\texttt{embeddings} extra), which upgrades the embedding fidelity
without changing the interface or downstream analysis.

The embeddings support semantic retrieval over the corpus and feed two
visualizations: a PCA two-dimensional projection ((Figure pca
embeddings)) that maps the topical geography of the literature, and a
hierarchical clustering dendrogram ((Figure dendrogram)) that reveals
the similarity structure of the document collection.
\end{frame}

\end{document}
